\section{Grids and chosen configurations}\label{app:grids}

The search spaces are as follows. LEAF-LGN and the other post-hoc backbones use learning rate $\{10^{-2}, 3\times10^{-3}\}$ and weight decay $\{10^{-4}, 10^{-5}\}$, with the backbone's own grid. LEAF-LGN uses $w\in\{0, 0.25, 0.5\}$ (its policy weight). EASE uses $w\in\{0,0.25,0.5\}$, $\lambda\in\{10^{3},3\times10^{3},10^{4}\}$ and $\beta\in\{0,0.25,0.5\}$, which gives 27 configurations. The DAX grid has 32 configurations: prior strength $\{80,160\}$, centre $\{$tag, content$\}$, lifelong strength $k\in\{0,100\}$, tab factor on or off, and stratum cells off or on with $\texttt{prior\_s}=80$. The gammas are $\{0, 0.05, 0.1, 0.15, 0.2, 0.3, 0.4, 0.5, 0.7, 1.0\}$. A post-hoc arm searches the backbone trainings times 32 DAX configurations times 10 gammas, for example $12\times32\times10=3{,}840$ for LEAF-LGN+DAX and $27\times32\times10=8{,}640$ for EASE+DAX. The QP family with DAX searches 10 backbone trainings, or $10\times32\times10=3{,}200$, and the 1K arm searches 20, or 6,400 (Appendix~\ref{app:qp}).

The tables below give the configuration chosen on validation in each split, for every method with a result file. Entries marked TBD have no result file yet. In the main run, LEAF-LGN+DAX chooses the tab factor on, the stratum cell prior 80 and gamma 0.5 in every split; the attractiveness centre is the tag mean in four splits and the content centre in split two; the lifelong strength is 100 in every split, the upper edge of its grid, and the prior strength is 80 in every split, the lower edge (Section~\ref{sec:limits}). The backbone chooses learning rate 3e-3, weight decay 1e-4 and $w\in\{0.25,0.5\}$, with 14 to 27 epochs.

\begin{table}[tbp]
\centering
\footnotesize
\caption{Chosen configurations per split: LEAF-LGN+DAX (main, post-hoc) (nDCG@10 mean 0.4854).}
\label{tab:grid00}
\begin{adjustbox}{max width=\linewidth}
\begin{tabular}{rp{0.82\linewidth}}
\toprule
Split & Chosen configuration \\
\midrule
1 & \texttt{lr=0.003, wd=0.0001, w=0.25, epoch=23, gamma=0.5, prior=80, centre=tag, k\_life=100, tab=True, prior\_s=80} \\
2 & \texttt{lr=0.003, wd=0.0001, w=0.5, epoch=25, gamma=0.5, prior=80, centre=content, k\_life=100, tab=True, prior\_s=80} \\
3 & \texttt{lr=0.003, wd=0.0001, w=0.5, epoch=27, gamma=0.5, prior=80, centre=tag, k\_life=100, tab=True, prior\_s=80} \\
4 & \texttt{lr=0.003, wd=0.0001, w=0.25, epoch=14, gamma=0.5, prior=80, centre=tag, k\_life=100, tab=True, prior\_s=80} \\
5 & \texttt{lr=0.003, wd=0.0001, w=0.5, epoch=27, gamma=0.5, prior=80, centre=tag, k\_life=100, tab=True, prior\_s=80} \\
\bottomrule
\end{tabular}
\end{adjustbox}
\end{table}

\begin{table}[tbp]
\centering
\footnotesize
\caption{Chosen configurations per split: LEAF-LGN (post-hoc base) (nDCG@10 mean 0.4736).}
\label{tab:grid01}
\begin{adjustbox}{max width=\linewidth}
\begin{tabular}{rp{0.82\linewidth}}
\toprule
Split & Chosen configuration \\
\midrule
1 & \texttt{lr=0.003, wd=0.0001, w=0, epoch=25} \\
2 & \texttt{lr=0.003, wd=0.0001, w=0, epoch=18} \\
3 & \texttt{lr=0.003, wd=0.0001, w=0, epoch=38} \\
4 & \texttt{lr=0.003, wd=1e-05, w=0, epoch=21} \\
5 & \texttt{lr=0.003, wd=0.0001, w=0, epoch=17} \\
\bottomrule
\end{tabular}
\end{adjustbox}
\end{table}

\begin{table}[tbp]
\centering
\footnotesize
\caption{Chosen configurations per split: QP family (stored, per-split best variant) (nDCG@10 mean 0.4634).}
\label{tab:grid02}
\begin{adjustbox}{max width=\linewidth}
\begin{tabular}{rp{0.82\linewidth}}
\toprule
Split & Chosen configuration \\
\midrule
1 & \texttt{lr=0.003, wd=0.0001, dim=64, layers=2, fn=1, epoch=38} \\
2 & \texttt{lr=0.01, wd=0.0001, dim=64, layers=2, fn=1, epoch=12} \\
3 & \texttt{lr=0.01, wd=0.0001, dim=64, layers=2, fn=1, epoch=11} \\
4 & \texttt{lr=0.003, wd=0.0001, dim=64, layers=2, fn=1, epoch=39} \\
5 & \texttt{lr=0.01, wd=1e-05, dim=64, layers=2, fn=1, epoch=3} \\
\bottomrule
\end{tabular}
\end{adjustbox}
\end{table}

\begin{table}[tbp]
\centering
\footnotesize
\caption{Chosen configurations per split: QP family+DAX (nDCG@10 mean 0.4733).}
\label{tab:grid03}
\begin{adjustbox}{max width=\linewidth}
\begin{tabular}{rp{0.82\linewidth}}
\toprule
Split & Chosen configuration \\
\midrule
1 & \texttt{method=DR-BPR-LGN-1, lr=0.003, wd=0.0001, dim=64, layers=2, fn=1, epoch=39, gamma=0.5, prior=80, centre=tag, k\_life=100, tab=True, prior\_s=None} \\
2 & \texttt{method=DR-BPR-LGN-1, lr=0.01, wd=0.0001, dim=64, layers=2, fn=1, epoch=12, gamma=0.5, prior=80, centre=tag, k\_life=100, tab=True, prior\_s=None} \\
3 & \texttt{method=DR-BPR-LGN-1, lr=0.003, wd=0.0001, dim=64, layers=2, fn=1, epoch=39, gamma=0.5, prior=80, centre=tag, k\_life=100, tab=True, prior\_s=None} \\
4 & \texttt{method=DR-BPR-LGN-1, lr=0.01, wd=0.0001, dim=64, layers=2, fn=1, epoch=19, gamma=0.5, prior=80, centre=tag, k\_life=100, tab=True, prior\_s=None} \\
5 & \texttt{method=DR-BPR-LGN-1, lr=0.01, wd=1e-05, dim=64, layers=2, fn=1, epoch=14, gamma=0.7, prior=80, centre=tag, k\_life=100, tab=True, prior\_s=None} \\
\bottomrule
\end{tabular}
\end{adjustbox}
\end{table}

\begin{table}[tbp]
\centering
\footnotesize
\caption{Chosen configurations per split: EASE (nDCG@10 mean 0.4507).}
\label{tab:grid04}
\begin{adjustbox}{max width=\linewidth}
\begin{tabular}{rp{0.82\linewidth}}
\toprule
Split & Chosen configuration \\
\midrule
1 & \texttt{w=0.5, epoch=0, lam=10000.0, beta=0.0} \\
2 & \texttt{w=0.5, epoch=0, lam=10000.0, beta=0.0} \\
3 & \texttt{w=0.5, epoch=0, lam=10000.0, beta=0.0} \\
4 & \texttt{w=0.5, epoch=0, lam=10000.0, beta=0.0} \\
5 & \texttt{w=0.5, epoch=0, lam=10000.0, beta=0.0} \\
\bottomrule
\end{tabular}
\end{adjustbox}
\end{table}

\begin{table}[tbp]
\centering
\footnotesize
\caption{Chosen configurations per split: EASE+DAX (nDCG@10 mean 0.4623).}
\label{tab:grid05}
\begin{adjustbox}{max width=\linewidth}
\begin{tabular}{rp{0.82\linewidth}}
\toprule
Split & Chosen configuration \\
\midrule
1 & \texttt{w=0.5, epoch=0, gamma=0.15, prior=80, centre=tag, k\_life=100, tab=True, prior\_s=None, lam=10000.0, beta=0.0} \\
2 & \texttt{w=0.5, epoch=0, gamma=0.15, prior=80, centre=tag, k\_life=100, tab=True, prior\_s=80, lam=10000.0, beta=0.0} \\
3 & \texttt{w=0.5, epoch=0, gamma=0.2, prior=80, centre=tag, k\_life=100, tab=True, prior\_s=None, lam=10000.0, beta=0.0} \\
4 & \texttt{w=0.5, epoch=0, gamma=0.15, prior=80, centre=tag, k\_life=100, tab=True, prior\_s=None, lam=10000.0, beta=0.0} \\
5 & \texttt{w=0.5, epoch=0, gamma=0.15, prior=80, centre=tag, k\_life=100, tab=True, prior\_s=None, lam=10000.0, beta=0.0} \\
\bottomrule
\end{tabular}
\end{adjustbox}
\end{table}

\begin{table}[tbp]
\centering
\footnotesize
\caption{Chosen configurations per split: LightGCN (nDCG@10 mean 0.4576).}
\label{tab:grid06}
\begin{adjustbox}{max width=\linewidth}
\begin{tabular}{rp{0.82\linewidth}}
\toprule
Split & Chosen configuration \\
\midrule
1 & \texttt{lr=0.003, wd=0.0001, dim=64, layers=2, epoch=38} \\
2 & \texttt{lr=0.01, wd=0.0001, dim=64, layers=2, epoch=14} \\
3 & \texttt{lr=0.01, wd=0.0001, dim=64, layers=2, epoch=6} \\
4 & \texttt{lr=0.01, wd=0.0001, dim=64, layers=2, epoch=12} \\
5 & \texttt{lr=0.01, wd=1e-05, dim=64, layers=2, epoch=3} \\
\bottomrule
\end{tabular}
\end{adjustbox}
\end{table}

\begin{table}[tbp]
\centering
\footnotesize
\caption{Chosen configurations per split: r-AdjNorm (nDCG@10 mean 0.4620).}
\label{tab:grid07}
\begin{adjustbox}{max width=\linewidth}
\begin{tabular}{rp{0.82\linewidth}}
\toprule
Split & Chosen configuration \\
\midrule
1 & \texttt{lr=0.01, wd=0.0001, dim=64, layers=2, r=0.7, epoch=17} \\
2 & \texttt{lr=0.01, wd=0.0001, dim=64, layers=2, r=0.7, epoch=15} \\
3 & \texttt{lr=0.01, wd=0.0001, dim=64, layers=2, r=0.7, epoch=11} \\
4 & \texttt{lr=0.01, wd=1e-05, dim=64, layers=2, r=0.7, epoch=18} \\
5 & \texttt{lr=0.01, wd=0.0001, dim=64, layers=2, r=0.7, epoch=12} \\
\bottomrule
\end{tabular}
\end{adjustbox}
\end{table}

\begin{table}[tbp]
\centering
\footnotesize
\caption{Chosen configurations per split: DIN (27K) (nDCG@10 mean 0.4675).}
\label{tab:grid08}
\begin{adjustbox}{max width=\linewidth}
\begin{tabular}{rp{0.82\linewidth}}
\toprule
Split & Chosen configuration \\
\midrule
1 & \texttt{lr=0.01, wd=1e-05, epoch=9} \\
2 & \texttt{lr=0.01, wd=1e-05, epoch=9} \\
3 & \texttt{lr=0.01, wd=0.0001, epoch=7} \\
4 & \texttt{lr=0.01, wd=0.0001, epoch=7} \\
5 & \texttt{lr=0.01, wd=1e-05, epoch=8} \\
\bottomrule
\end{tabular}
\end{adjustbox}
\end{table}

\begin{table}[tbp]
\centering
\footnotesize
\caption{Chosen configurations per split: DIN+DAX (27K) (nDCG@10 mean 0.4796).}
\label{tab:grid09}
\begin{adjustbox}{max width=\linewidth}
\begin{tabular}{rp{0.82\linewidth}}
\toprule
Split & Chosen configuration \\
\midrule
1 & \texttt{lr=0.01, wd=1e-05, epoch=9, gamma=0.4, prior=80, centre=tag, k\_life=100, tab=True, prior\_s=None} \\
2 & \texttt{lr=0.01, wd=1e-05, epoch=9, gamma=0.4, prior=80, centre=tag, k\_life=100, tab=True, prior\_s=80} \\
3 & \texttt{lr=0.003, wd=0.0001, epoch=10, gamma=0.4, prior=80, centre=tag, k\_life=100, tab=True, prior\_s=None} \\
4 & \texttt{lr=0.01, wd=0.0001, epoch=7, gamma=0.4, prior=80, centre=tag, k\_life=100, tab=True, prior\_s=None} \\
5 & \texttt{lr=0.01, wd=1e-05, epoch=9, gamma=0.4, prior=80, centre=tag, k\_life=100, tab=True, prior\_s=80} \\
\bottomrule
\end{tabular}
\end{adjustbox}
\end{table}

\begin{table}[tbp]
\centering
\footnotesize
\caption{Chosen configurations, SIM (27K): not available.}
\label{tab:grid10}
\begin{adjustbox}{max width=\linewidth}
\begin{tabular}{l}
\toprule
Split \\
\midrule
\textnormal{\{TBD:SIM (27K) chosen configurations\}} \\
\bottomrule
\end{tabular}
\end{adjustbox}
\end{table}

\begin{table}[tbp]
\centering
\footnotesize
\caption{Chosen configurations, SIM+DAX (27K): not available.}
\label{tab:grid11}
\begin{adjustbox}{max width=\linewidth}
\begin{tabular}{l}
\toprule
Split \\
\midrule
\textnormal{\{TBD:SIM+DAX (27K) chosen configurations\}} \\
\bottomrule
\end{tabular}
\end{adjustbox}
\end{table}

\begin{table}[tbp]
\centering
\footnotesize
\caption{Chosen configurations per split: LEAF-LGN+DAX (1K) (nDCG@10 mean 0.4590).}
\label{tab:grid12}
\begin{adjustbox}{max width=\linewidth}
\begin{tabular}{rp{0.82\linewidth}}
\toprule
Split & Chosen configuration \\
\midrule
1 & \texttt{lr=0.01, wd=0.0001, w=0.25, epoch=0, gamma=0.3, prior=160, centre=tag, k\_life=0, tab=False, prior\_s=80} \\
2 & \texttt{lr=0.01, wd=0.0001, w=0.5, epoch=0, gamma=0.5, prior=80, centre=tag, k\_life=100, tab=False, prior\_s=None} \\
3 & \texttt{lr=0.01, wd=0.0001, w=0.5, epoch=0, gamma=0.5, prior=160, centre=tag, k\_life=100, tab=False, prior\_s=80} \\
4 & \texttt{lr=0.01, wd=0.0001, w=0.5, epoch=3, gamma=0.2, prior=80, centre=content, k\_life=100, tab=True, prior\_s=None} \\
5 & \texttt{lr=0.01, wd=0.0001, w=0.5, epoch=1, gamma=0.15, prior=80, centre=content, k\_life=100, tab=False, prior\_s=80} \\
\bottomrule
\end{tabular}
\end{adjustbox}
\end{table}

\begin{table}[tbp]
\centering
\footnotesize
\caption{Chosen configurations per split: QP family+DAX (1K) (nDCG@10 mean 0.4358).}
\label{tab:grid13}
\begin{adjustbox}{max width=\linewidth}
\begin{tabular}{rp{0.82\linewidth}}
\toprule
Split & Chosen configuration \\
\midrule
1 & \texttt{method=DR-BPR-LGN-1, lr=0.01, wd=0.0001, dim=64, layers=2, fn=1, epoch=4, gamma=1, prior=160, centre=tag, k\_life=100, tab=True, prior\_s=80} \\
2 & \texttt{method=DR-BPR-LGN-p, lr=0.01, wd=0.0001, dim=64, layers=2, fn=1, pbeta=0.5, epoch=37, gamma=1, prior=160, centre=tag, k\_life=100, tab=True, prior\_s=80} \\
3 & \texttt{method=DR-BPR-LGN-p, lr=0.01, wd=1e-05, dim=64, layers=2, fn=1, pbeta=0.5, epoch=31, gamma=0.4, prior=80, centre=tag, k\_life=100, tab=False, prior\_s=80} \\
4 & \texttt{method=DR-BPR-LGN-p, lr=0.01, wd=0.0001, dim=64, layers=2, fn=1, pbeta=0.5, epoch=33, gamma=0.7, prior=80, centre=tag, k\_life=100, tab=False, prior\_s=80} \\
5 & \texttt{method=DR-BPR-LGN-1, lr=0.01, wd=0.0001, dim=64, layers=2, fn=1, epoch=10, gamma=0.7, prior=80, centre=tag, k\_life=100, tab=True, prior\_s=None} \\
\bottomrule
\end{tabular}
\end{adjustbox}
\end{table}

\begin{table}[tbp]
\centering
\footnotesize
\caption{Chosen configurations per split: DIN (1K) (nDCG@10 mean 0.4189).}
\label{tab:grid14}
\begin{adjustbox}{max width=\linewidth}
\begin{tabular}{rp{0.82\linewidth}}
\toprule
Split & Chosen configuration \\
\midrule
1 & \texttt{lr=0.003, wd=0.0001, epoch=0} \\
2 & \texttt{lr=0.003, wd=0.0001, epoch=0} \\
3 & \texttt{lr=0.003, wd=0.0001, epoch=0} \\
4 & \texttt{lr=0.01, wd=1e-05, epoch=1} \\
5 & \texttt{lr=0.003, wd=0.0001, epoch=0} \\
\bottomrule
\end{tabular}
\end{adjustbox}
\end{table}

\begin{table}[tbp]
\centering
\footnotesize
\caption{Chosen configurations per split: DIN+DAX (1K) (nDCG@10 mean 0.4228).}
\label{tab:grid15}
\begin{adjustbox}{max width=\linewidth}
\begin{tabular}{rp{0.82\linewidth}}
\toprule
Split & Chosen configuration \\
\midrule
1 & \texttt{lr=0.003, wd=0.0001, epoch=0, gamma=0.7, prior=160, centre=tag, k\_life=100, tab=False, prior\_s=None} \\
2 & \texttt{lr=0.01, wd=0.0001, epoch=1, gamma=1, prior=160, centre=tag, k\_life=100, tab=True, prior\_s=80} \\
3 & \texttt{lr=0.003, wd=0.0001, epoch=0, gamma=0.7, prior=160, centre=tag, k\_life=100, tab=False, prior\_s=None} \\
4 & \texttt{lr=0.01, wd=1e-05, epoch=1, gamma=0.4, prior=80, centre=tag, k\_life=100, tab=False, prior\_s=80} \\
5 & \texttt{lr=0.003, wd=0.0001, epoch=1, gamma=1, prior=160, centre=tag, k\_life=100, tab=True, prior\_s=None} \\
\bottomrule
\end{tabular}
\end{adjustbox}
\end{table}

\begin{table}[tbp]
\centering
\footnotesize
\caption{Chosen configurations per split: SIM (1K) (nDCG@10 mean 0.4246).}
\label{tab:grid16}
\begin{adjustbox}{max width=\linewidth}
\begin{tabular}{rp{0.82\linewidth}}
\toprule
Split & Chosen configuration \\
\midrule
1 & \texttt{lr=0.003, wd=0.0001, epoch=0} \\
2 & \texttt{lr=0.01, wd=0.0001, epoch=1} \\
3 & \texttt{lr=0.003, wd=0.0001, epoch=0} \\
4 & \texttt{lr=0.003, wd=1e-05, epoch=7} \\
5 & \texttt{lr=0.003, wd=0.0001, epoch=1} \\
\bottomrule
\end{tabular}
\end{adjustbox}
\end{table}

\begin{table}[tbp]
\centering
\footnotesize
\caption{Chosen configurations per split: SIM+DAX (1K) (nDCG@10 mean 0.4357).}
\label{tab:grid17}
\begin{adjustbox}{max width=\linewidth}
\begin{tabular}{rp{0.82\linewidth}}
\toprule
Split & Chosen configuration \\
\midrule
1 & \texttt{lr=0.003, wd=0.0001, epoch=0, gamma=0.2, prior=160, centre=content, k\_life=100, tab=False, prior\_s=None} \\
2 & \texttt{lr=0.01, wd=1e-05, epoch=1, gamma=0.4, prior=160, centre=tag, k\_life=0, tab=True, prior\_s=80} \\
3 & \texttt{lr=0.01, wd=1e-05, epoch=1, gamma=0.3, prior=160, centre=tag, k\_life=100, tab=True, prior\_s=None} \\
4 & \texttt{lr=0.003, wd=0.0001, epoch=6, gamma=0.5, prior=80, centre=tag, k\_life=0, tab=False, prior\_s=None} \\
5 & \texttt{lr=0.01, wd=1e-05, epoch=1, gamma=0.5, prior=160, centre=tag, k\_life=100, tab=True, prior\_s=80} \\
\bottomrule
\end{tabular}
\end{adjustbox}
\end{table}

